\begin{center}
    {\LARGE \textbf{2025无机化学与物理化学}} \\[2ex]
    {\large 时量180分钟 \quad 满分150分}
\end{center}

\vspace{0.5cm}
\noindent\textbf{注意事项:}
\begin{enumerate}[label=\arabic*., parsep=0pt, itemsep=0pt]
    \item 答题前,考生先将自己的姓名、学号、考场号、座位号填写在答题卡上,将条形码准确粘贴在考生信息条形码粘贴区。
    \item 考试结束后,将本试卷与答题纸一并收回。
    \item 回答各个问题时,将答案写在答题纸上,写在本试卷上无效。
\end{enumerate}

\vspace{1cm}
\begin{center}
    {\Large \textbf{无机化学部分(共75分)}}
\end{center}

\vspace{0.5cm}
\noindent\textbf{一、选择题（每题2分，共20分）}

\begin{enumerate}[label=\arabic*., itemsep=1.5ex]
    \item 下列关于过渡金属化合物的颜色为无色的是
    \begin{tasks}(4)
        \task \ce{CrF3}
        \task \ce{MnF3}
        \task \ce{TiCl3}
        \task \ce{ScF3}
    \end{tasks}

    \item 下列选项稳定性顺序正确的是
    \begin{tasks}(1)
        \task \ce{[Co(NH3)6]^3+} $>$ \ce{[Co(SCN)4]^2-} $>$ \ce{[Co(CN)6]^3-}
        \task \ce{[HgI4]^2-} $>$ \ce{[HgCl4]^2-} $>$ \ce{[Hg(CN)4]^2-}
        \task \ce{[Fe(SCN)6]^3-} $>$ \ce{[Fe(CN)6]^3-} $>$ \ce{[Fe(SCN)4]+}
        \task \ce{[Ni(en)3]^2+} $>$ \ce{[Ni(NH3)6]^2+} $>$ \ce{[Ni(H2O)6]^2+}
    \end{tasks}

    \item 下列物质不含离域 $\pi$ 键的是
    \begin{tasks}(4)
        \task \ce{ClO2}
        \task \ce{SOCl2}
        \task \ce{HN3}
        \task \ce{SO3}
    \end{tasks}

    \item 下列物质有3种几何异构体的是
    \begin{tasks}(2)
        \task \ce{[Pt(NH3)2BrCl3]}
        \task \ce{[Co(NH3)2(en)Cl2]+}
        \task \ce{[PtCl2(NO2)2(NH3)2]}
        \task \ce{[Co(NH3)(en)Cl3]}
    \end{tasks}

    \item 合成氨工艺中吸收 $\ce{H2}$ 中的杂质 $\ce{CO}$ 的是
    \begin{tasks}(2)
        \task \ce{[Cu(NH3)2]Ac}
        \task \ce{[Ag(NH3)2]OH}
        \task \ce{[Cu(NH3)4](OH)2}
        \task \ce{[Cu(NH3)4](Ac)2}
    \end{tasks}

    \item 第四周期第ⅤA族的元素中，符合量子数 $m=0$ 的电子数有几个
    \begin{tasks}(4)
        \task 22
        \task 21
        \task 15
        \task 14
    \end{tasks}

    \item 下列选项错误的是
    \begin{tasks}(1)
        \task \ce{I2} 不能氧化 \ce{Fe^2+}，但在 \ce{KCN} 存在下，\ce{I2} 可以氧化 \ce{Fe^2+}
        \task \ce{Cu} 可以溶于 \ce{KCN} 并置换出 \ce{H2}，而 \ce{Ag} 在通入空气中才可以溶于 \ce{KCN}
        \task 赤血盐可由 \ce{Fe2(SO4)3} 和 \ce{KCN} 直接在水溶液中制备
        \task \ce{CdS} 可以溶于 $3\ \mathrm{mol\cdot L^{-1}}$ \ce{HCl}，但不溶于 $3\ \mathrm{mol\cdot L^{-1}}$ \ce{HClO4}
    \end{tasks}

    \item 下列可以用 \ce{NaOH} 分离的是
    \begin{tasks}(2)
        \task \ce{Ag+} 和 \ce{Hg^2+}
        \task \ce{Hg^2+} 和 \ce{Zn^2+}
        \task \ce{Pb^2+} 和 \ce{Cr^3+}
        \task \ce{Cr^3+} 和 \ce{Al^3+}
    \end{tasks}

    \item 下列选项正确的是
    \begin{tasks}(1)
        \task 半径：$\ce{Co} > \ce{Ni}$，$\ce{Ni} > \ce{Cu}$
        \task 第一电离能：$\ce{Fe} > \ce{Ru}$，$\ce{Ru} < \ce{Os}$
        \task 第一电子亲和能：$\ce{B} > \ce{C}$，$\ce{C} < \ce{Si}$
        \task 电负性：$\ce{O} < \ce{Cl}$，$\ce{O} < \ce{F}$
    \end{tasks}

    \item 下列选项错误的是
    \begin{tasks}(1)
        \task \ce{CoCl2 * 6H2O} 受热脱水不生成 \ce{HCl}
        \task \ce{Pb2O3}、\ce{Co2O3} 和 \ce{Ni2O3} 与浓盐酸作用均有氯气放出
        \task \ce{NaHCO3}、\ce{CaCO3} 和 \ce{Na2CO3} 的热稳定性依次增大
        \task \ce{(NH4)2SO4}、\ce{NH4NO2} 和 \ce{NH4NO3} 受热分解时均无 \ce{NH3} 放出
    \end{tasks}
\end{enumerate}

\vspace{0.5cm}
\noindent\textbf{二、简要回答下列问题。本题共7小题，第1题15分，2~7题每题5分，共45分。}

\begin{enumerate}[label=\arabic*., itemsep=1.5ex]
    \item 解释或给出实验现象并写出必要的相关化学反应方程式：
    \begin{enumerate}[label=(\arabic*), noitemsep, topsep=0pt, partopsep=0pt]
        \item 将 \ce{Zn} 单质投入酸性重铬酸钾溶液，溶液由橙色经绿色而变为蓝色，久置后又变成绿色；
        \item 向 \ce{Fe2(SO4)3} 溶液中依次加入饱和 \ce{(NH4)2C2O4} 溶液、少量 \ce{KSCN} 溶液、少量盐酸和过量 \ce{NH4F} 溶液；
        \item 将 \ce{MnO2}、\ce{KClO3} 和 \ce{KOH} 固体混合后用煤气灯加热，得到绿色固体。用大量水处理绿色固体得到紫色溶液和棕黑色沉淀，再通入过量 \ce{NO2} 则沉淀和颜色消失，得到无色溶液。
    \end{enumerate}

    \item 试比较 \ce{(CH3)3N} 和 \ce{(SiH3)3N} 的碱性和空间结构，并解释原因。

    \item 以方铅矿(\ce{PbS})为原料制备铅单质和铅丹。

    \item 指出 \ce{[Co(NO2)(NH3)5]Cl2} 和 \ce{[Co(ONO)(NH3)5]Cl2} 中哪个是黄色的，哪个是红色的，并说明原因。

    \item 用 \ce{Zn(OH)2} 的氨水溶液可以鉴别黑色固体 \ce{FeS} 和 \ce{Fe2S3}，试给出反应现象并解释原因。

    \item 向 \ce{Bi^3+} 和 \ce{Sn^2+} 混合澄清溶液中加入 \ce{NaOH}，有黑色沉淀生成。给出黑色沉淀是什么，并解释为何 $\text{pH}$ 的改变导致该反应的发生。

    \item 结合化学方程式描述马氏试砷法的化学依据，并指出试样是否含砷的关键实验现象。
\end{enumerate}

\vspace{0.5cm}
\noindent\textbf{三、推断题：本题共1小题，共10分。}

白色固体化合物 A 不溶于水，加热后剧烈分解，生成固体 B 和气体 C。固体 B 不溶于水、盐酸，但溶于热的稀硝酸得到无色溶液 D 和无色气体 E，E 暴露在空气中迅速转化为另一种棕色气体。向溶液 D 中滴加盐酸，得到白色固体 F，F 不溶于过量浓盐酸。气体 C 不与一般试剂作用，但金属镁能在 C 气氛下燃烧，得到黄色固体 G。化合物 A 以 \ce{H2S} 溶液处理时得到黑色沉淀 H、无色溶液 I 和气体 C。过滤后，固体 H 溶于浓硝酸得到气体 E、黄色沉淀 J 和溶液 D。用 \ce{HCl} 溶液处理滤液 I 得到臭鸡蛋味气体，用 \ce{NaOH} 溶液处理滤液 I 得到可以使湿润的红色石蕊试纸变蓝色的无色刺激性气体。

据以上，推断字母 A $\sim$ J 所代表的物质。

\vspace{1cm}
\begin{center}
    {\Large \textbf{物理化学部分(共75分)}}
\end{center}